\documentclass[11pt,a4paper]{article}
\usepackage[utf8]{inputenc}\usepackage[T1]{fontenc}\usepackage[spanish]{babel}
\usepackage[margin=2cm]{geometry}\usepackage{amsmath,amssymb,booktabs,hyperref}
\hypersetup{colorlinks=true,urlcolor=blue}\setlength{\parindent}{0pt}\setlength{\parskip}{6pt}
\begin{document}
\begin{center}{\Large Límites conjuntos del escalamiento por filas}\\[3pt]
Nota preliminar independiente, 15 de octubre de 2026\end{center}
\textbf{Estado.} Derivación y certificados comprobados; novedad editorial y
ventaja práctica no establecidas. OMS v1.3.5 permanece sin cambios. Ninguna
llamada nueva a optimizadores.

\textbf{Antecedente clásico.} Para $A$ cuadrada no singular,
$S(A)=\|\,|A^{-1}|\,|A|\,\|_\infty$ es invariante por filas y
$\min_D\kappa_\infty(DA)=S(A)$; equilibrar las sumas absolutas de filas alcanza
el mínimo. Véase \href{https://eprints.maths.manchester.ac.uk/651/1/SURVEY.pdf}{Higham (1994), ecuaciones 2.11--2.12}.
No es un resultado nuevo de estas notas. El escalamiento óptimo tiene extensa
literatura, incluyendo \href{https://www.tuhh.de/ti3/paper/rump/Ru03a.pdf}{Rump (2003)}
y \href{https://arxiv.org/abs/2209.00809}{Qu et al.} Una búsqueda preliminar no
determina prioridad para la extensión acotada siguiente.

\textbf{Intervalos.} Sean $r_i=\sum_j|a_{ij}|$, $C=|A^{-1}|$ y
$0<L_i\leq\alpha_i\leq H_i<\infty$. Definir
\[
 t_0=\max_i L_i r_i,\quad
 \alpha_i^*=\min\{H_i,t_0/r_i\},\quad
 K_* = \max_j\sum_i C_{ji}\max\{r_i,t_0/H_i\}.
\]
Entonces $\alpha^*$ es admisible y alcanza el mínimo $K_*$ de
$\kappa_\infty(D_\alpha A)$. Es una derivación corta propia aquí, sin afirmar
novedad. Seleccionar factores no exige la inversa; certificar $K_*$ en general sí.

\emph{Prueba.} Para cualquier $\alpha$ admisible, $t=\max_i\alpha_i r_i\geq t_0$,
\[
 \kappa_\infty(D_\alpha A)=\max_j\sum_i C_{ji}\frac{t}{\alpha_i},\qquad
 \frac{t}{\alpha_i}\geq\max\{r_i,t/H_i\}\geq\max\{r_i,t_0/H_i\}.
\]
Esto prueba la cota inferior; $\alpha^*$ tiene norma máxima de filas $t_0$
y alcanza todas esas cotas, probando optimalidad.

\textbf{Qué detecta.} Intervalos locales no vacíos no garantizan
$\kappa_\infty\leq K$ si $K<K_*$. Para $A=I$, cajas $[1,10^9]$ y
$[1/16,1/8]$ dan mínimo $8$, aun con $S(A)=1$. La matriz
$\left[\begin{smallmatrix}1&1\\1&1+2^{-20}\end{smallmatrix}\right]$ tiene
piso $4\,194\,307$, aunque casi no hay dispersión de coeficientes dentro de filas.
Son ejemplos construidos, no casos operativos del despacho.

\newpage
\textbf{Compatibilidad conjunta.} Para $Ax=b$, la hipótesis de residuo
exportado $|D(A\widehat x-b)|\leq\epsilon\mathbf1$ da
$|\widehat x-x|\leq\epsilon C(1/\alpha)$. Es la mejor envolvente sobre la caja
independiente de residuos; otras restricciones pueden excluir sus extremos.
No se deduce esa hipótesis de un parámetro del solver.

Exigir a la vez $\kappa_\infty(DA)\leq K$ y esa envolvente $\leq\Delta$ equivale
a la factibilidad del LP en $(v,t)$:
\[
 v_i\geq r_i,\qquad t/H_i\leq v_i\leq t/L_i,\qquad
 Cv\leq K\mathbf1,\qquad \epsilon Cv\leq t\Delta.
\]
Las primeras restricciones implican $t\geq t_0>0$. Necesidad: usar
$t=\max_i\alpha_i r_i$, $v_i=t/\alpha_i$. Suficiencia: tomar
$\alpha_i=t/v_i$; la norma de filas es $\leq t$ y las dos últimas filas
garantizan los objetivos. Es una transformación lineal fraccional elemental.

\textbf{Certificado exacto.} Para
\[
 A=\begin{pmatrix}1&-1\\0&1\end{pmatrix},\quad
 L=(1,1),\ H=(1,128),\ \epsilon=1,\ K=10,\ \Delta=(2,1/128),
\]
condición sola es factible con $\alpha=(1,2)$, dando $\kappa=3$.
Error solo es factible con $(1,128)$, pero da $\kappa=129$.
La caja y la segunda precisión fuerzan esa última pareja, así que los objetivos
son incompatibles juntos. Todos esos factores son potencias de dos.
El LP tiene el certificado Farkas (multiplicadores no negativos)
\[
\begin{aligned}
129(-v_2\leq-1)&+(v_1+v_2\leq10)\\
&+(t-v_1\leq0)+128(v_2-t/128\leq0)
\end{aligned}
\quad\Longrightarrow\quad 0\leq-119.
\]
No usa tolerancias ni etiquetas de un optimizador. Relajar $K$ a $129$ admite
$(v_1,v_2,t)=(128,1,128)$. $\epsilon=1$ es una normalización del ejemplo.

\textbf{Grilla dyádica.} Reducir la caja a extremos dyádicos admisibles,
resolver el mínimo continuo y redondear cada factor hacia abajo conserva la
caja. La norma de filas no aumenta y la norma de la inversa crece por menos de
dos. Se obtiene una aproximación por factor dos al óptimo dyádico de condición.
Esto no preserva automáticamente una meta conjunta $\Delta$ ni garantiza
representabilidad del vector solución. La inviabilidad del LP real sí descarta
factores dyádicos. La representabilidad de productos debe entrar en los extremos.

\newpage
\textbf{Bloque hidráulico real.} En tres MILP originales de 360 períodos,
con 1440 enteras intactas, se extraen las columnas de volumen en sus balances,
fijando flujos y vertimientos. El bloque $B$ tiene diagonal $1$ y subdiagonal
$-1$; $B^{-1}$ tiene unos en el triángulo inferior. Se comprueban sus entradas
y $BB^{-1}=I$. Por eso $\kappa_\infty(B)=720$ y el mínimo por filas es $719$.
Lo alcanza $D=\operatorname{diag}(1,1/2,\ldots,1/2)$, admisible para los tres
escenarios y con exportación exacta. Los prefijos analíticos $T=8,24$ tienen
pisos $15,47$; no son instancias MILP ni resultados LP--DW.

\begin{center}\begin{tabular}{lrrr}\toprule
Caso & Inicial & Mínimo & Error acumulado máximo, m$^3$\\\midrule
001 & 720 & 719 & 0.0022245 \\
022 & 720 & 719 & 3.3755e-06 \\
046 & 720 & 719 & 0.0035253 \\
\bottomrule\end{tabular}\end{center}

Se auditan los 1093 eventos del paquete anterior, incluyendo finales repetidos.
Se suma exactamente el residuo de agua hasta cada período y se verifica el
telescopado directo. Es discrepancia respecto de reconstruir almacenamiento
con los mismos flujos/inicialización; no un error físico medido ni reparación
del MILP. Aunque 100 m$^3$ por balance permitiría una cota genérica de 36\,000
m$^3$ al final, el máximo observado es menor que 0.004 m$^3$. No aparece una
necesidad práctica en este banco.

La cota terminal teórica al usar factores máximos es aproximadamente
$1.296\cdot10^{-9}$ m$^3$, siempre condicionada a un residuo exportado que
realmente se cumpla. No es residuo observado ni promesa del solver. Los
escenarios autorizados no exigen esa precisión extrema.

\textbf{Límites y decisión.} El piso 719 corresponde a ese subsistema con flujos
fijos: no es una cota de condición de la matriz rectangular completa, de todas
las bases ni de la matriz interna del solver. La norma infinito no es la norma
1/2 del resto de los diagnósticos. Cambiar las holguras unitarias al construir
una base puede dar una matriz distinta de $DB$. La caracterización requiere
un sistema cuadrado fijo y puede necesitar una inversa densa costosa.

Se verifican 4667 comparaciones exactas de grillas en 24 matrices pequeñas,
el certificado Farkas y cuatro tests. Esto comprueba las fórmulas, no prioridad
histórica ni generalización. La parte libre es clásica; la formulación acotada
y conjunta requiere revisar más antecedentes y demostrar utilidad real antes
de considerar OMS. No se añade un teorema al artículo por sí solo.

README.md describe el replay sin solvers. La auditoría completa usa el ZIP padre
fijado por SHA256; no se incluye otra copia de sus 30 MB. Los nueve paquetes
previos y la entrega OMS permanecen sin cambios.
\end{document}
